% 第4章: インターフェースで統一する
\section{インターフェースで統一する}

%% ── Q1: モジュールが増えたら ──────────────────
\begin{qabox}{モジュールが増えたらどうするの？}

LED、温度センサー、モーター……とモジュールが増えていくと、こんなコードになりがちです。

\begin{badexample}{呼び方がバラバラ}
\begin{lstlisting}
void setup() {
    led.init();
    tempSensor.begin();       // beginだったり
    motor.start();            // startだったり
    display.initialize();     // initializeだったり
}

void loop() {
    led.toggle();
    tempSensor.readValue();   // メソッド名もバラバラ
    motor.drive();
    display.refresh();
}
\end{lstlisting}
\end{badexample}

モジュールごとに初期化メソッドの名前が違う、更新メソッドの名前が違う。これでは新しいモジュールを追加するたびに「このモジュールは何を呼べばいいんだっけ？」と調べる羽目になります。

もっと困るのは、「全モジュールの初期化を順番に実行する」ような共通処理が書けないことです。

\end{qabox}

%% ── Q2: IModuleの役割 ──────────────────────
\begin{qabox}{IModuleって何をしてくれるの？}

\texttt{IModule}は「すべてのモジュールが守るべきルール」を定めたインターフェースです。

\begin{goodexample}{IModuleインターフェース}
\begin{lstlisting}
struct SystemData;  // 前方宣言（次の章で詳しく説明）

class IModule {
public:
    bool enabled = true;

    virtual bool init() = 0;
    virtual void updateInput(SystemData& data) {}
    virtual void updateOutput(SystemData& data) {}
    virtual void deinit() {}
    virtual ~IModule() {}
};
\end{lstlisting}
\end{goodexample}

ここで出てくる \texttt{virtual} と \texttt{= 0} について説明しましょう。

\subsection{virtualとは}

\texttt{virtual}は「この関数は子クラスで上書き（オーバーライド）できますよ」という印です。

\subsection{= 0（純粋仮想関数）とは}

\texttt{= 0}は「この関数は子クラスで必ず実装してね。実装しないとコンパイルエラーにするよ」という意味です。IModuleでは \texttt{init()} だけが \texttt{= 0} なので、必ず実装が必要です。

一方、\texttt{updateInput()} や \texttt{updateOutput()} は \texttt{\{\}}（空の中括弧）で「何もしないデフォルト実装」が書かれています。入力モジュールは \texttt{updateInput()} を、出力モジュールは \texttt{updateOutput()} をオーバーライドします。両方を持つモジュール（BLE等）は両方をオーバーライドします。

\subsection{配列で回せるメリット}

すべてのモジュールがIModuleを継承すると、こんなことができます。

\begin{goodexample}{配列でまとめて初期化}
\begin{lstlisting}
// 異なるクラスでも、IModule*の配列に入れられる
IModule* modules[] = {
    &led,
    &tempSensor,
    &motor,
};

// まとめて初期化
for (int i = 0; i < 3; i++) {
    if (!modules[i]->init()) {
        Serial.println("init failed!");
        modules[i]->enabled = false;
    }
}
\end{lstlisting}
\end{goodexample}

更新は入力モジュールと出力モジュールで配列を分け、それぞれ \texttt{updateInput()} と \texttt{updateOutput()} を呼びます。この仕組みは第5章「3フェーズモデル」で詳しく解説します。

\texttt{LedModule}も\texttt{TempSensorModule}も\texttt{MotorModule}も、すべて\texttt{IModule*}として統一的に扱えます。新しいモジュールを追加しても、for文はそのままです。

\begin{notebox}[overrideキーワード]
子クラスで親の仮想関数を実装するときは、\texttt{override}キーワードを付けましょう。タイプミスで別の関数を定義してしまうバグを防げます。
\begin{lstlisting}
class LedModule : public IModule {
public:
    bool init() override { /* ... */ }      // OK
    void updateOutput(SystemData& data) override; // OK
};
\end{lstlisting}
\end{notebox}

\end{qabox}

%% ── Q3: enabledの意味 ──────────────────────
\begin{qabox}{なんでenabledがあるの？}

ハードウェアは壊れることがあります。センサーが接触不良で応答しない、ディスプレイの初期化に失敗した……そんなとき、プログラム全体を止めたくはありません。

\texttt{enabled}フラグを使えば、初期化に失敗したモジュールだけをスキップできます。

\begin{lstlisting}
for (int i = 0; i < MODULE_COUNT; i++) {
    if (!modules[i]->init()) {
        Serial.println("init failed, disabling");
        modules[i]->enabled = false;  // このモジュールだけ無効化
    }
}

void loop() {
    for (int i = 0; i < INPUT_COUNT; i++) {
        if (inputModules[i]->enabled) {    // 無効なモジュールはスキップ
            inputModules[i]->updateInput(systemData);
        }
    }
    // ... 出力モジュールも同様にenabledチェック
}
\end{lstlisting}

温度センサーが壊れても、LEDやモーターは動き続けます。組み込みシステムでは「部分的に動く」ことが重要です。

\end{qabox}
